\section{Conclusion}
\label{sec:conclusion}

This paper has described the PolicyEngine tax--benefit microsimulation as the PolicyEngine Macro suite uses it: statute encoded as a graph of dated, parameterised variables, resolved at household resolution for the United Kingdom and the United States, driven through an integration layer that validates the reform object, records provenance, and returns a score in a schema shared with the suite's macro members.

Its central claim is a boundary. A household calculation is exact arithmetic---no sampling, no weights, no imputation, no estimated coefficients---and can be reproduced with a pencil, which is why the registry records its uncertainty as ``none.'' A population estimate is a weighted sum of that same exact arithmetic over an \emph{estimated} population, and inherits sampling, measurement, calibration and ageing error from the survey beneath it, which is why the registry records survey and calibration uncertainty in the same sentence and after a semicolon. The two are different in kind, not in degree, and the practical consequence is asymmetric: improving the rule engine barely moves a population costing, while improving the microdata does.

The validation posture follows the same split. There is no forecast error to report because nothing is forecast; the rules are checked against implemented legislation, and one worked case---\pounds 50{,}000 of employment income yielding \pounds 7{,}486 of income tax, \pounds 2{,}994 of National Insurance and \pounds 39{,}520 of net income in 2026---is offered as a check the reader performs rather than a claim the reader accepts. The single committed external benchmark, 1p on the UK basic rate at \pounds 6.46bn in 2026 rising to \pounds 7.38bn by 2030 against HMRC's \pounds 6.9bn rising to \pounds 8.2bn, comes in \emph{below} the official figure at every year that can be matched---by 6.4 per cent in the first scored year, widening thereafter---for reasons that are understood and are not defects to be tuned away: administrative versus survey data, a post-behavioural official basis against a static one, threshold-freeze assumptions. It is one reform, one country, one tax, one direction, measured against another estimate of a different object, with only its first year both year-matched and gated at all; the paper declines to call it a validation.

What the model cannot do is short and absolute: GDP, inflation, interest rates, macro feedback. The suite supplies those through bridges, and the paper has been as explicit about the bridges' limits as about the model's. The route into the OBR emulator discards the distribution the microsimulation computed and injects a single scalar per quarter as a held add-factor on nominal household disposable income; it is a demand-side approximation, not a general model of reform incidence, and its output should be read as sign and order of magnitude.

The clearest piece of outstanding work is the one this paper's central section had to admit. The uncertainty on a population estimate is documented as a sentence and not as a number: no interval, no standard error, no propagated imputation variance. A distinction between exact arithmetic and estimated aggregates is worth more when the second half has a range attached to it, and computing design-consistent variance from the survey's own replicate weights is the natural first step. Beyond that: adopting the newer FRS vintage, building a committed US benchmark to match the UK one, checking distributional output against published analyses rather than only for internal consistency, and carrying incidence rather than a scalar across the macro bridge.
